\chapter{Programar com IA sem Perder o Controlo}
\label{ch:software}

Um programa recebe três medições: zero, quatro e um valor em falta. Uma proposta calcula a média de quatro. Corre sem erro e devolve um número plausível. Perdeu o zero, que era uma medição válida. Antes de escolher uma biblioteca ou redesenhar o programa, convém escrever a regra que faltou.

\section{Escolher o tamanho da ajuda}

Podes pedir a um assistente que complete uma função, que explique uma alteração ou que compare duas arquiteturas. No primeiro caso, dá a assinatura e exemplos. No segundo, mostra a chamada que importa e o comportamento observado. No terceiro, indica as restrições e a decisão que está em aberto. Um pedido de correção para a média começa nos três valores, na regra e no resultado esperado.

O termo \emph{vibe coding} descreve usos em que a pessoa orienta a geração por intenção e aceita propostas com pouca leitura. Pode servir para explorar um protótipo. Se o código for usado por outras pessoas, prepara uma revisão e explica como verificar o comportamento que lhes interessa.

\begin{promptlabbox}{Identificar o modo certo antes de pedir ajuda}
\textbf{Objetivo e contexto:} Delimitar um pedido de programação. Usa um exemplo teu autorizado ou o cálculo da média deste capítulo.

Formula um pedido para uma função e outro para rever a integração dessa função num relatório. Indica o que pode mudar e que comportamento os consumidores já esperam.

\textbf{Resultado esperado e validação:} O primeiro pedido permite escrever a função e as verificações. O segundo inclui a origem dos dados, os consumidores e o tratamento de falhas. Se pedires propostas ao modelo, confirma se ficam no âmbito escolhido.

\textbf{Falhas possíveis:} Misturar correção e redesenho, aceitar código para uma biblioteca diferente ou omitir o comportamento esperado para dados antigos.
\end{promptlabbox}

\section{Ler o código a partir de uma chamada}

Escolhe um caminho concreto. Para a média, pede que o assistente siga \texttt{[0, 4, None]}, indique o ramo escolhido para cada elemento e mostre o resultado. Confronta a explicação com o código. Numa aplicação maior, pergunta quem chama a função, que entradas lhe chegam, que estado modifica e que exceção atravessa a fronteira. Um diagrama pode ajudar a seguir as chamadas, desde que cada ligação corresponda ao fluxo que verificaste.

\begin{lstlisting}
Segue a chamada com a entrada indicada.
Localiza onde cada valor é aceite ou descartado.
Indica o resultado esperado a partir da regra.
Se propuseres uma alteração, identifica primeiro
os chamadores e os comportamentos que preserva.
\end{lstlisting}

\section{Uma regra independente da implementação}

Neste exemplo didático, a função recebe inteiros pequenos ou \texttt{None}. Ignora só os valores em falta. Se não houver medições, devolve \texttt{None}. Para $[0,4]$, a conta é $(0+4)/2=2$. Uma filtragem pela veracidade do valor descarta zero. O caso foi escolhido para separar as duas interpretações.

\begin{lstlisting}
def media(valores):
    presentes = [v for v in valores if v is not None]
    if not presentes:
        return None
    return sum(presentes) / len(presentes)

assert media([0, 4, None]) == 2
assert media([0]) == 0
assert media([]) is None
assert media([None]) is None
\end{lstlisting}

As quatro verificações foram executadas localmente. O resultado esperado vem da regra e da conta, que chamamos \emph{oráculo} do teste. Se o teste partir da mesma proposta que descartou zero, é fácil codificar quatro como resultado esperado e obter um sucesso enganador.

\begin{promptlabbox}{Refinar um pedido de geração de código}
\textbf{Objetivo e contexto:} Comparar uma descrição curta com uma especificação que cobre entradas, resultados e fronteiras.

Escreve primeiro «faz uma função que calcule a média». Numa sessão separada, dá a regra acima com o domínio de entrada. Reserva o caso com zero para a verificação. Examina as duas propostas antes de as executar num ambiente isolado.

\textbf{Resultado esperado e validação:} Cada versão é avaliada pelo mesmo oráculo. Inclui lista vazia, valor em falta e zero válido. Regista quando ambas produzem o resultado correto.

\textbf{Falhas possíveis:} Aceitar testes que repetem a implementação, alterar o domínio entre comparações ou tratar uma exceção inesperada como comportamento aprovado.
\end{promptlabbox}

\section{Da função ao repositório}

Numa alteração real, identifica o ponto de entrada, o ambiente e as versões das dependências. Trabalha numa branch e executa os testes com dados de teste. Antes de aplicar uma sugestão, revê a diferença completa, incluindo comandos, ficheiros, permissões e operações externas. Para uma proposta de otimização, começa por reproduzir o problema e mede com entradas representativas. A função acima percorre a sequência e cria uma lista de valores presentes. Se essa lista ficar cara, mede o efeito antes de alterar a implementação.

Pede triagem de injeções, validação de entradas e exposição de dados onde o código cruza uma fronteira. Confirma a ocorrência no fluxo e consulta a documentação da versão usada. Quando uma proposta se parece com código de outra origem, verifica a proveniência e as condições de utilização antes de a incorporar.

Na descrição da mudança, separa o que esperavas, o que observaste e a dúvida funcional que sobra. Uma sugestão pode passar nos testes da função e falhar quando um consumidor interpreta \texttt{None} como zero. O capítulo~\ref{ch:spec-first} amplia essa conversa para a especificação do sistema.
